% GENERATED FILE. Edit canonical lessons, then run npm run generate:books.
% canonical-source-hash: fnv1a64:607499c6b6fdf9fa
% canonical-lessons: MR-C455-dhasti, MR-C455-becain, MR-C455-dagdag, MR-C455-talmalne, MR-C455-ghusmatne

\chapter{Dread and Restless Nights}
\label{ch:mr-dread-and-restless-nights}
\hlchaptermodality{455}

\begin{chapteropening}
\textbf{By the end of this chapter:} \emph{I can talk about dread, a hectic rush and a restless night.}

Complete the last lesson of chapter 455: I can talk about dread, a hectic rush and a restless night.
\end{chapteropening}

\section[\texorpdfstring{\mr{धास्ती} (dhāstī)}{धास्ती (dhāstī)}]{\mr{धास्ती} (dhāstī) — dread}

\label{lesson:MR-C455-dhasti}



\begin{quote}
Before the new one: say the Marathi for to coax round, then the Marathi for to be crestfallen.
\end{quote}

\subsection*{You'll want to know: \mr{धास्ती}}
\textbf{\mr{धास्ती}} — \emph{dhāstī} — \textquotedblleft{}dread\textquotedblright{}.

\begin{quote}
\mr{परीक्षेची} \mr{धास्ती} \mr{तिच्या} \mr{मनात} \mr{होती}.
\end{quote}

\begin{quote}
\emph{parīkṣecī dhāstī ticyā manāt hotī.}
\end{quote}

\begin{quote}
Dread of the exam was on her mind.
\end{quote}

Feminine.

Feminine: \mr{परीक्षेची} \mr{धास्ती}, dread of the exam. \mr{धास्ती} \mr{वाटणे}, to feel dread, with the dative.

\begin{cousinweb}
\textbf{\mr{धास्ती}} (\emph{dhāstī}) is a Marathi noun whose origin is not settled.
\end{cousinweb}

\begin{grammarlens}[title={What you've built}]
One more word for describing what you lived and felt.
\end{grammarlens}

\subsection*{Your turn}
\begin{itemize}
\raggedright
  \item \emph{Say it:} \emph{dhāstī}
  \item \emph{Say it:} \emph{dhāstī}, once more
  \item \emph{Say it:} \emph{parīkṣecī dhāstī ticyā manāt hotī.} — Dread of the exam was on her mind.
  \item \emph{Recall:} say the Marathi for to coax round, then the Marathi for to be crestfallen, then say \emph{dhāstī} again
\end{itemize}

\begin{tcolorbox}[breakable,colback=teal!4,colframe=teal!35!black,title={Before you move on}]
What does \mr{धास्ती} mean? (\textquotedblleft{}Dread\textquotedblright{}.) Say it once more.
\end{tcolorbox}
\section[\texorpdfstring{\mr{बेचैन} (becain)}{बेचैन (becain)}]{\mr{बेचैन} (becain) — restless; uneasy}

\label{lesson:MR-C455-becain}



\begin{quote}
Before the new one: say the Marathi for to be crestfallen, then the Marathi for dread.
\end{quote}

\subsection*{You'll want to know: \mr{बेचैन}}
\textbf{\mr{बेचैन}} — \emph{becain} — \textquotedblleft{}restless; uneasy\textquotedblright{}.

\begin{quote}
\mr{बाळ} \mr{रात्रभर} \mr{बेचैन} \mr{होतं}.
\end{quote}

\begin{quote}
\emph{bāḷ rātrabhar becain hotaṁ.}
\end{quote}

\begin{quote}
The baby was restless all night.
\end{quote}

Unchanging. The prefix \mr{बे}- means without: here, without rest.

\begin{cousinweb}
\textbf{\mr{बेचैन}} (\emph{becain}) is a loan from Persian \emph{bechain}, \textquotedblleft{}restless\textquotedblright{}, built on \emph{be-}, \textquotedblleft{}without\textquotedblright{}, and \emph{chain}, \textquotedblleft{}rest\textquotedblright{}.
\end{cousinweb}

\begin{grammarlens}[title={What you've built}]
One more word for describing what you lived and felt.
\end{grammarlens}

\subsection*{Your turn}
\begin{itemize}
\raggedright
  \item \emph{Say it:} \emph{becain}
  \item \emph{Say it:} \emph{becain}, once more
  \item \emph{Say it:} \emph{bāḷ rātrabhar becain hotaṁ.} — The baby was restless all night.
  \item \emph{Recall:} say the Marathi for to be crestfallen, then the Marathi for dread, then say \emph{becain} again
\end{itemize}

\begin{tcolorbox}[breakable,colback=teal!4,colframe=teal!35!black,title={Before you move on}]
What does \mr{बेचैन} mean? (\textquotedblleft{}Restless; uneasy\textquotedblright{}.) Say it once more.
\end{tcolorbox}
\section[\texorpdfstring{\mr{दगदग} (dagdag)}{दगदग (dagdag)}]{\mr{दगदग} (dagdag) — a wearing rush; exhausting bustle}

\label{lesson:MR-C455-dagdag}



\begin{quote}
Before the new one: say the Marathi for dread, then the Marathi for restless.
\end{quote}

\subsection*{You'll want to know: \mr{दगदग}}
\textbf{\mr{दगदग}} — \emph{dagdag} — \textquotedblleft{}a wearing rush; exhausting bustle\textquotedblright{}.

\begin{quote}
\mr{रोजच्या} \mr{दगदगीमुळे} \mr{आई} \mr{थकते}.
\end{quote}

\begin{quote}
\emph{rojcyā dagdagīmuḷe āī thakte.}
\end{quote}

\begin{quote}
Mother gets tired from the daily rush.
\end{quote}

Feminine.

Feminine; oblique \mr{दगदगी}-. \mr{दगदग} \mr{करणे}, to wear oneself out rushing about.

\begin{cousinweb}
\textbf{\mr{दगदग}} (\emph{dagdag}) is a Marathi noun built by doubling a syllable; its origin is not settled.
\end{cousinweb}

\begin{grammarlens}[title={What you've built}]
One more word for describing what you lived and felt.
\end{grammarlens}

\subsection*{Your turn}
\begin{itemize}
\raggedright
  \item \emph{Say it:} \emph{dagdag}
  \item \emph{Say it:} \emph{dagdag}, once more
  \item \emph{Say it:} \emph{rojcyā dagdagīmuḷe āī thakte.} — Mother gets tired from the daily rush.
  \item \emph{Recall:} say the Marathi for dread, then the Marathi for restless, then say \emph{dagdag} again
\end{itemize}

\begin{tcolorbox}[breakable,colback=teal!4,colframe=teal!35!black,title={Before you move on}]
What does \mr{दगदग} mean? (\textquotedblleft{}A wearing rush; exhausting bustle\textquotedblright{}.) Say it once more.
\end{tcolorbox}
\section[\texorpdfstring{\mr{तळमळणे} (taḷmaḷṇe)}{तळमळणे (taḷmaḷṇe)}]{\mr{तळमळणे} (taḷmaḷṇe) — to toss and turn; to fret}

\label{lesson:MR-C455-talmalne}



\begin{quote}
Before the new one: say the Marathi for restless, then the Marathi for a wearing rush.
\end{quote}

\subsection*{You'll want to know: \mr{तळमळणे}}
\textbf{\mr{तळमळणे}} — \emph{taḷmaḷṇe} — \textquotedblleft{}to toss and turn; to fret\textquotedblright{}.

\begin{quote}
\mr{काळजीमुळे} \mr{तो} \mr{रात्रभर} \mr{तळमळला}.
\end{quote}

\begin{quote}
\emph{kāḷjīmuḷe to rātrabhar taḷmaḷlā.}
\end{quote}

\begin{quote}
He tossed and turned all night with worry.
\end{quote}

Past (he …): \mr{तो} \mr{तळमळला}, \emph{to taḷmaḷlā}.

Of a sleepless night or a restless longing; \mr{तळमळ}, the restlessness itself.

\begin{cousinweb}
\textbf{\mr{तळमळणे}} (\emph{taḷmaḷṇe}) is a Marathi verb built on a doubled syllable, \mr{तळ}-\mr{मळ}; its origin is not settled.
\end{cousinweb}

\begin{grammarlens}[title={What you've built}]
One more word for describing what you lived and felt.
\end{grammarlens}

\subsection*{Your turn}
\begin{itemize}
\raggedright
  \item \emph{Say it:} \emph{taḷmaḷṇe}
  \item \emph{Say it:} \emph{taḷmaḷṇe}, once more
  \item \emph{Say it:} \emph{kāḷjīmuḷe to rātrabhar taḷmaḷlā.} — He tossed and turned all night with worry.
  \item \emph{Recall:} say the Marathi for restless, then the Marathi for a wearing rush, then say \emph{taḷmaḷṇe} again
\end{itemize}

\begin{tcolorbox}[breakable,colback=teal!4,colframe=teal!35!black,title={Before you move on}]
What does \mr{तळमळणे} mean? (\textquotedblleft{}To toss and turn; to fret\textquotedblright{}.) Say it once more.
\end{tcolorbox}
\section[\texorpdfstring{\mr{घुसमटणे} (ghusmaṭṇe)}{घुसमटणे (ghusmaṭṇe)}]{\mr{घुसमटणे} (ghusmaṭṇe) — to feel stifled}

\label{lesson:MR-C455-ghusmatne}



\begin{quote}
Before the new one: say the Marathi for a wearing rush, then the Marathi for to toss and turn.
\end{quote}

\subsection*{You'll want to know: \mr{घुसमटणे}}
\textbf{\mr{घुसमटणे}} — \emph{ghusmaṭṇe} — \textquotedblleft{}to feel stifled\textquotedblright{}.

\begin{quote}
\mr{गर्दीच्या} \mr{बसमध्ये} \mr{मी} \mr{घुसमटलो}.
\end{quote}

\begin{quote}
\emph{gardīcyā basmadhye mī ghusmaṭlo.}
\end{quote}

\begin{quote}
I felt stifled in the crowded bus.
\end{quote}

Past (he …): \mr{तो} \mr{घुसमटला}, \emph{to ghusmaṭlā}.

Of a body in a stuffy room, and of a mind with no one to talk to: \mr{तिचा} \mr{जीव} \mr{घुसमटत} \mr{होता}.

\begin{cousinweb}
\textbf{\mr{घुसमटणे}} (\emph{ghusmaṭṇe}) is a Marathi verb whose origin is not settled.
\end{cousinweb}

\begin{grammarlens}[title={What you've built}]
One more word for describing what you lived and felt.
\end{grammarlens}

\subsection*{Your turn}
\begin{itemize}
\raggedright
  \item \emph{Say it:} \emph{ghusmaṭṇe}
  \item \emph{Say it:} \emph{ghusmaṭṇe}, once more
  \item \emph{Say it:} \emph{gardīcyā basmadhye mī ghusmaṭlo.} — I felt stifled in the crowded bus.
  \item \emph{Recall:} say the Marathi for a wearing rush, then the Marathi for to toss and turn, then say \emph{ghusmaṭṇe} again
\end{itemize}

\begin{tcolorbox}[breakable,colback=teal!4,colframe=teal!35!black,title={Before you move on}]
What does \mr{घुसमटणे} mean? (\textquotedblleft{}To feel stifled\textquotedblright{}.) Say it once more.
\end{tcolorbox}
